
\documentclass[manuscript=article]{achemso}
\usepackage{booktabs}
\usepackage{amsmath}
\author{A. Placeholder}
\affiliation{Department of Verified Templates, Placeholder Institute, Placeholder City, Country}
\email{verified@example.org}
\author{B. Example}
\affiliation{Journal Systems Lab, Example University, Example City, Country}
\title{ACS Sensors: Official ACS Submission Preview}
\begin{document}
\begin{abstract}
This submission-style preview is rendered with the official ACS achemso class in manuscript article mode. The content remains placeholder text, but the article is intentionally structured to expose chemistry-oriented front matter, section flow, figure balance, table treatment, and references in a form that is closer to a practical ACS submission.
\end{abstract}
\section{Introduction}
The introduction is written as compact scientific prose so the preview reads like a real ACS manuscript rather than a bare template check page.
\section{Related Work}
The related-work section is explicit so the page rhythm reveals how prior chemistry literature, methods, and experiments separate within the manuscript.
\section{Methodology}
The methodology section introduces representative notation and a figure placeholder so the ACS template can be inspected under more realistic technical content.
\begin{equation}
\mathcal{S} = \sum_{i=1}^{M} w_i \left\| y_i - \hat{y}_i \right\|_2^2.
\end{equation}
\begin{figure}[ht]
\centering
\fbox{\parbox[c][2.0in][c]{0.82\linewidth}{\centering ACS Figure Placeholder\\[0.4em]\normalsize Sensor schematic, assay pipeline, or molecular workflow}}
\caption{Representative ACS figure block used to inspect manuscript-mode spacing, caption treatment, and the transition from methods text to the first major figure.}
\end{figure}
\section{Experiments}
The experiments section provides enough structure to expose table handling, subsection rhythm, and the transition from methods to quantitative comparison.
\subsection{Experimental Setup}
This paragraph acts as a placeholder for experimental conditions, sensing platform details, sample preparation, and baseline methods.
\subsection{Quantitative Results}
\begin{table}[ht]
\caption{Quantitative comparison block used to inspect the official ACS manuscript template under realistic submission-style content.}
\centering
\begin{tabular}{lcc}
\toprule
Method & Score & Time \\
\midrule
Placeholder A & 0.91 & 1.4 \\
Placeholder B & 0.88 & 1.7 \\
Placeholder C & 0.84 & 2.0 \\
\bottomrule
\end{tabular}
\end{table}
\section{Conclusion and Discussion}
This page verifies the official ACS template path while preserving the six-part manuscript skeleton requested for the library. The result is intended to read more like a review-ready chemistry submission than a simple placeholder sheet.
\begin{thebibliography}{9}
\bibitem{ref1} A. Smith and B. Lee. Template-aware ACS manuscript preparation. \textit{ACS Placeholder Journal} \textbf{2025}, \textit{1}, 1--10.
\bibitem{ref2} C. Garcia and D. Patel. Figure and table balance in chemistry manuscripts. \textit{ACS Sensors} \textbf{2026}, \textit{2}, 20--31.
\end{thebibliography}
\end{document}
